\documentclass[11pt]{article}
\usepackage[margin=1in]{geometry}
\usepackage[T1]{fontenc}
\usepackage{lmodern}
\usepackage{amsmath,amssymb,booktabs,array,longtable}
\usepackage{listings,xcolor,hyperref}
\setlength{\emergencystretch}{8em}
\makeatletter
\newcommand{\sha}{\begingroup\catcode`\_=12 \sha@read}
\newcommand{\sha@read}[1]{\ttfamily\sha@loop#1\relax\endgroup}
\def\sha@loop#1{%
  \ifx\relax#1\else
    #1\allowbreak
    \expandafter\sha@loop
  \fi}
\makeatother
\hypersetup{colorlinks=true,linkcolor=blue,urlcolor=blue,citecolor=blue,
  pdftitle={A replication and formalization of the Mi three-trait QTL SEM},pdfauthor={}}
\newcommand{\proofcommit}{05d69ea141adf3d5ee936d3b48061bde241250dc}
\newcommand{\proofbase}{https://github.com/alectionary/vibe-math/blob/\proofcommit/formalization/correlated-residuals}
\newcommand{\lean}[2]{\href{\proofbase/#1.lean\#L#2}{\texttt{#2}}}
\lstset{basicstyle=\footnotesize\ttfamily,columns=fullflexible,keepspaces=true,
  breaklines=true,showstringspaces=false,frame=single,framerule=.3pt,
  literate={∀}{{$\forall$}}1 {∃}{{$\exists$}}1 {≠}{{$\ne$}}1 {≤}{{$\le$}}1
  {≥}{{$\ge$}}1 {→}{{$\to$}}1 {←}{{$\leftarrow$}}1 {ℕ}{{$\mathbb{N}$}}1
  {ℝ}{{$\mathbb{R}$}}1 {θ}{{$\theta$}}1 {μ}{{$\mu$}}1 {ω}{{$\omega$}}1
  {φ}{{$\varphi$}}1 {γ}{{$\gamma$}}1
  {Φ}{{$\Phi$}}1 {Ψ}{{$\Psi$}}1 {Γ}{{$\Gamma$}}1 {Σ}{{$\Sigma$}}1
  {Ω}{{$\Omega$}}1 {∫}{{$\int$}}1 {∂}{{$\partial$}}1 {·}{{$\cdot$}}1
  {⟨}{{$\langle$}}1 {⟩}{{$\rangle$}}1
  {ᵀ}{{$^{\mathsf{T}}$}}1 {ᵐ}{{$^{m}$}}1 {¹}{{$^{1}$}}1
  {⁻}{{$^{-}$}}1 {₀}{{$_{0}$}}1 {₁}{{$_{1}$}}1 {₂}{{$_{2}$}}1 {₃}{{$_{3}$}}1
  {∧}{{$\land$}}1 {∨}{{$\lor$}}1}

\title{A replication and formalization of the three-trait QTL SEM\\
in Mi (2009) and Mi, Eskridge, Wang, and coauthors (2010)}
\author{}
\date{27 October 2026}

\newcommand{\ev}[1]{\textup{\emph{(#1)}}}

\begin{document}
\maketitle

\noindent
This manuscript was drafted by a language model as a replication and
formalization record. The date is an internal manuscript date. The note is
not a journal article, and it does not speak for the authors of the sources
it studies.

\begin{abstract}
This note records an attempt by a language model to replicate, and to
formalize in Lean, the three-trait QTL structural equation model associated
with Mi's 2009 University of Nebraska--Lincoln dissertation, with
Mi, Eskridge, Wang, Baenziger, Campbell, Gill, Dweikat, and Bovaird,
Statistical Applications in Genetics and Molecular Biology~9 (2010),
article~38, and with the 2008 conference version of that design. The
dissertation was not obtained. On the pattern printed in equations~(7)--(8)
of the 2010 paper, diagonal residuals are a left inverse of the covariance,
each single residual covariance collides, and each of the three bow-free
patterns is a left inverse. Those statements are Lean proofs, with no
\texttt{sorry}. A first reading of equation~(7), with one minus in front of
$M$, does not reproduce Table~2's total effect $0.75$. Measuring the glyph
positions puts the minuses on individual entries. That reading reproduces
the parameter block, including $0.75$, and raises the single-trait power for
the third QTL on the first trait from $13\%$ to $86\%$ in $1000$ replicates.
The printed multi-trait powers $64$, $77$, and $96$ are not jointly
reproduced by any design this note could reconcile with the simulation
paragraph. The closest percentage triple found is $65.2$, $73.3$, $99.3$.
A product of binomial probabilities for $64$, $77$, and $96$ at the
baseline rates is about $1.4\times 10^{-12}$, under an independence
assumption the three detections do not satisfy. Two further comparisons are
left as source discrepancies under our reading: Table~2 prints $-0.435$ and
$0.065$ where the arithmetic gives $-0.4375$ and $0.0625$, and Table~1
prints SEM powers above the least-squares powers on a model where the two
likelihood ratios coincide.
\end{abstract}

\section*{Question, source, and answer}

\begin{center}
\small
\begin{tabular}{@{}p{0.22\textwidth}p{0.72\textwidth}@{}}
\toprule
Question & Can the trait residuals in the Mi et al.\ (2010) QTL SEM be correlated after environment and block effects are removed? Printed page 18 (PDF page 20). \\
Source & Mi et al., SAGMB 9 (2010), article 38. The request is not a numbered conjecture. \\
Answer on the printed pattern & On equation (7), diagonal residuals are a left inverse of the covariance. Each single residual covariance collides. Each bow-free missing-edge pattern is a left inverse, and the plug-in is almost surely consistent. Theorems \lean{CorrelatedResiduals/Diagonal}{246}, \lean{CorrelatedResiduals/Collision}{114}, \lean{CorrelatedResiduals/BowFree}{187}, and \lean{CorrelatedResiduals/BowConsistency}{270}. \\
What the simulations add & A seeded replication, the bias and size distortion of the diagonal fit, and coverage of \texttt{recoverBow}. These are sampled experiments, not Lean proofs. \\
Still open & A fit to the wheat measurements, the general mixed-graph characterization, and a Lean central limit theorem. \\
\bottomrule
\end{tabular}
\end{center}

The manuscript date 27 October 2026 is an internal date. It is not an external publication date. Elapsed times below are measured to that date. A year-only source is reported as a range.

\begin{center}
\small
\begin{tabular}{@{}p{0.28\textwidth}p{0.22\textwidth}p{0.22\textwidth}p{0.20\textwidth}@{}}
\toprule
Source & Earliest date used & Inspected version & Elapsed to 27 Oct 2026 \\
\midrule
Brito and Pearl, bow-free identification & 2002 (year) & SEM 9, 459--474 & 23--24 years \\
Campbell et al., chromosome 3A QTLs & 2003 (year) & Crop Sci.\ 43, 1493--1505 & 22--23 years \\
Li et al., structural QTL model & 21 July 2006 & PLoS Genet.\ 2, e114 & 20 years, 3 months, 6 days \\
Mi et al., the source question & 2010 (year) & SAGMB 9, article 38 & 15 years 10 months to 16 years 10 months \\
Valente et al., recursive search & June 2010 (month) & Genetics 185, 633--644 & 16 years, 4 months (month precision) \\
Drton, Foygel, and Sullivant & arXiv 4 Mar 2010; online 13 May 2011 & Ann.\ Statist.\ 39, 865--886 & 15 years, 5 months, 14 days from the journal date \\
Foygel, Draisma, and Drton, half-trek & arXiv 27 July 2011; online 3 Oct 2012 & Ann.\ Statist.\ 40, 1682--1713 & 14 years, 24 days from the journal date \\
Hidalgo-Contreras et al. & 28 September 2021 & Agronomy 11, 1953 & 5 years, 29 days \\
This manuscript & 27 October 2026 & internal manuscript date & --- \\
\bottomrule
\end{tabular}
\end{center}

\section{What this replication did}
\label{sec:did}

The work in this directory is a formalization of one identification question
and a set of seeded simulations aimed at Table~1 and Table~2 of Mi et al.\
(2010). The sections after this one keep the working record: the quoted
request, the glyph coordinates, the Lean statements, the design grid, and
the bow-free and Campbell calculations. Commands and file paths are collected
in Appendix~\ref{app:artifacts}.

\subsection{The Lean formalization}

The Lean development is
\texttt{formalization/correlated-residuals/CorrelatedResiduals/}, project
file \texttt{lakefile.toml}, toolchain Lean~4.22.0 and Mathlib v4.22.0.
\texttt{lake build} is the build. The proofs use no \texttt{sorry} and no
\texttt{sorryAx}. \texttt{\#print axioms} on the theorems named below reports
\texttt{propext}, \texttt{Classical.choice}, and \texttt{Quot.sound}.
\texttt{Classical.choice} enters through the matrix inverse and the
positive-definite unit argument. \texttt{Quot.sound} is the construction of
$\mathbb{R}$. \ev{our own Lean proof}

The statements, with file and line, are:
\begin{itemize}
\item \texttt{recover\_sigma} (\texttt{Diagonal.lean}:246). On the
equation-(7) zero pattern, positive diagonal $\Psi$ and positive definite
$\Phi$ make $(B,\Gamma,\Psi,\Phi)$ a function of $\Sigma$.
\item \texttt{bow12\_same\_law} (\texttt{Collision.lean}:114), and the two
sibling collisions for the other pairs. Each one-edge residual extension
collides.
\item \texttt{recover\_bow} (\texttt{BowFree.lean}:187),
\texttt{recover\_bow01} (\texttt{BowMissing.lean}:148), and
\texttt{recover\_bow23} (\texttt{BowMissing.lean}:353). Each bow-free
missing-edge pattern is a left inverse.
\item \texttt{unitUpper\_factor\_unique} (\texttt{Triangular.lean}:68) and
\texttt{sigmaRec\_injective} (\texttt{Triangular.lean}:298).
\item \texttt{plug\_in\_consistent\_ae} (\texttt{Law.lean}:174) and
\texttt{plug\_in\_bow\_consistent\_ae}
(\texttt{BowConsistency.lean}:270), with the two further bow patterns.
Almost-sure consistency of the plug-in. This is not a Lean central limit
theorem.
\end{itemize}
The Gaussian delta method is a written derivative and a numeric check
(\texttt{delta\_method.py}), not a Lean proof. The listings in
Appendix~\ref{app:lean} are the theorem bodies at commit
\sha{05d69ea141adf3d5ee936d3b48061bde241250dc}.

\subsection{The simulation studies}

Three seeded studies live in \texttt{simulation/studies.py}. Study~1 is a
100-replicate scan under the earlier reading of equation~(7). Study~2 is the
bow-free Monte Carlo, 400 replicates, seed \texttt{20101028}: the diagonal
plug-in sends the missing-edge coefficient to $0.8\rho$, and
\texttt{recoverBow} interval coverage at the nominal $95\%$ level falls
between about $0.915$ and $0.958$. Study~3 refits $B$ and $\Psi$ to three
phenotypic correlations from Campbell et al.\ (2003) and reruns the scan.
That run is a scale check. It is not a replication of Table~1, and the
Campbell PDF itself was not stored.
\ev{our own simulation}

A separate grid, \texttt{simulation/gaps.py}, reruns the scan under one
change at a time, including 1000-replicate versions of the earlier reading
and the glyph reading. \texttt{simulation/assumptions.py} crosses the
choices the simulation paragraph does not pin down, at 1000 replicates, and
also writes a noncentral chi-square approximation. \texttt{qtl1\_search.py}
moves the first QTL, the sample size, the marker gap, and the cofactor rule,
again at 1000 replicates, and records whether each choice is consistent with
the printed sentence.

\subsection{The replication results}

Under the glyph reading, $n=150$, seed \texttt{20101101}, and 1000
replicates, the wide-window SEM powers are $50.8$, $53.5$, and $86.1$.
A second seed, \texttt{20101102}, gives $49.4$, $48.3$, and $84.6$.
Single-trait powers for $Y_1$ are $3.2$, $1.2$, and $88.6$, against the
printed $6$, $3$, and $97$. Table~1 prints SEM powers $64$, $77$, and $96$,
least-squares powers $57$, $69$, and $94$, single-trait powers
$6/3/97$, $29/43/33$, and $10/1/17$, and overall powers $40$, $45$, and $97$.
\ev{our own simulation, set next to a cited table}

The glyph total-effect matrix is
\[
(I-M)^{-1}G
=\begin{pmatrix}1/8&0&3/4\\ 3/8&-7/16&3/8\\ 1/4&-1/8&1/4\end{pmatrix}.
\]
The $(1,3)$ entry is the $0.75$ in Table~2. The same arithmetic gives
$-7/16=-0.4375$ where Table~2 prints $-0.435$, and the indirect piece
$1/16=0.0625$ where the table prints $0.065$. Those two cells are discussed
with the unresolved comparisons. The $1000$-replicate single-trait power for
QTL3 on $Y_1$ is $86.2\%$ (seed \texttt{20101041}). The earlier reading,
which produced the total effect $0.25$, gave $13\%$ for that cell.

No configuration that stays with the printed positions, the printed sample
size, and the printed threshold reproduces $64$, $77$, and $96$ together.
The smallest squared error among the percentage triples that were computed
is $65.2$, $73.3$, $99.3$, from QTL1 at $24$~cM and $n=250$. That design is
not the design in the simulation sentence. Section~\ref{sec:open} records
what was tried and what is still open.

\section{Assumptions}
\label{sec:assumptions}

Each subsection is one assumption used to turn the printed paragraph into a
program. A choice the paragraph states is not listed here as an assumption.
The grid that turns one assumption off at a time is
Section~\ref{sec:gaps}.

\subsection{Design}

The population is an F1 selfed for twelve generations, with residual
heterozygotes coded $\pm 1$ at equal probability. Mi et al.\ (2010) say
recombinant inbred lines and do not give the generation count or a
segregation-distortion model. The map function in the SAGMB scan is Haldane.
The 2008 conference note names Haldane and prints the infinite-RIL expansion
$2r/(1+2r)$; the SAGMB simulation paragraph does not. Cofactors are every
marker except the flanking pair, which is the sentence on PDF page~10.
The conditional expectation of the QTL genotype given the flanking markers
follows Haley and Knott (1992), as the methods section says, implemented as
the exact two-point expectation on the twelve-generation pedigree.
\ev{assumption, flagged}

\subsection{Map positions}

Markers are placed at $0,10,\ldots,100$~cM. The sentence says eleven evenly
spaced markers on a $100$~cM chromosome. Putting a marker at each end is a
reading of ``evenly spaced,'' not a coordinate the paper prints. QTLs are
placed at $22$, $42$, and $78$~cM, which the sentence does print. The scan
step is $1$~cM. A detection is counted in a window of $\pm 10$~cM around the
true position, and the position standard deviations use that window. The
paper's position standard deviations are consistent with a window nearer
$\pm 5$~cM, which the text does not state. The $\pm 10$~cM window is ours.
\ev{assumption, flagged}

\subsection{Sample size}

The SAGMB simulation uses $150$ lines and $100$ replicates. The baseline
powers quoted above use $1000$ replicates of that $150$-line design, so that
Monte Carlo error is smaller than the gap under discussion. The $100$-replicate
size is what Table~1 itself is. Companion designs that were simulated and
then set aside are $250$ lines (Mi et al., Genetics Research, 2010) and
$300$ lines (the 2008 conference note). Neither is the SAGMB sample size.
\ev{assumption, flagged where the replicate count exceeds the printed $100$}

\subsection{Thresholds}

A multi-trait detection uses $\chi^2_{0.005,4}=14.86$. A single-trait
detection uses $\chi^2_{0.005,2}=10.60$. Both numbers are printed. Treating
$14.86$ as a LOD score, or as a $\chi^2$ critical value on $3$ or on $1$
degree of freedom, was tested and is not an assumption of the baseline.
The permutation threshold, a $95\%$ quantile near $16.1$ under the baseline
population, was computed for comparison and is not the threshold in the
scan that produced $50.8$, $53.5$, and $86.1$.
\ev{the printed thresholds, used as printed}

\subsection{Readings of equation~(7)}

Two readings were programmed.

The earlier reading puts the displayed minus in front of $M$ and takes every
entry of $G$ as positive, so $B=-M$. The total effect of QTL3 on $Y_1$ is
then $1/4$, and the single-trait power for that cell is $13\%$ at
$100$ replicates. This reading is study~1 of \texttt{studies.py} and the
rows of \texttt{gaps.py} labelled as the previous reading.

The glyph reading is $y=My+Gx+\varepsilon$, with
\[
M=\begin{pmatrix}0&1/2&1/4\\0&0&-1/2\\0&0&0\end{pmatrix},\qquad
G=\begin{pmatrix}-1/8&1/4&1/2\\1/2&-1/2&1/2\\1/4&-1/8&1/4\end{pmatrix},
\]
so $B=M$, and $\Psi=\mathrm{diag}(1.6,1.8,2.5)$ as printed in equation~(8).
The baseline, the joint search, and the Table~2 arithmetic use this reading.
How the signs were assigned is the next section.
\ev{our reading of the glyph positions}

\section{Ambiguities that were resolved}
\label{sec:resolved}

\subsection{The sign in front of \(M\)}

The first programs read equation~(7) as $y=-My+Gx+\varepsilon$. That
population does not produce Table~2's total effect $0.75$, and it produces
a single-trait power of $13\%$ where Table~1 prints $97\%$ for QTL3 on
$Y_1$.

The PDF content stream on page~10 of the ARS copy of Mi et al.\ (2010)
settles the operator. The equals sign ends at $x=145.9$ points. The
parenthesis of $M$ starts at $x=148.1$. No minus sits between them. Four
unary minuses are present, and each ends on the first digit of the number it
signs: $G_{11}=-0.125$ (the minus ends at $x=250.7$), $M_{23}=-0.5$ (ends at
$x=194.5$), $G_{22}=-0.5$ (ends at $x=316.7$), and $G_{32}=-0.125$ (ends at
$x=316.7$). \ev{our reading of the glyph positions}

The measurement is stored in
\texttt{simulation/out/qtl1/glyphs.json}, committed in
\sha{b4e8a3c594f12f419560d5cb2a569b649ae93a55}.
The command that re-reads the same page and checks those coordinates is
\begin{quote}\noindent
\sha{python3 formalization/correlated-residuals/simulation/measure_glyphs.py --pdf local-extraction/literature/unl-sem/mi-qtl2010.pdf --check formalization/correlated-residuals/simulation/out/qtl1/glyphs.json}
\end{quote}
The PDF is not in the repository. The copy used here has SHA-256
\sha{27c4f3926c9b5af9b399c9b3ef36918c28e7967ef777444b620279dd91628a81}.

\subsection{Table~2, including \(0.75\)}

With $B=M$ and the signed $G$ from the glyph reading,
$(I-B)^{-1}G$ has $(1,3)$ entry $3/4$. The rest of the parameter block
matches Table~2 except the two QTL2 cells discussed in
Section~\ref{sec:rounding}. The command is
\begin{quote}\noindent
\sha{python3 formalization/correlated-residuals/simulation/table2_totals.py}
\end{quote}
It writes \texttt{simulation/out/table2\_totals.json} and exits with status
$1$ if the $(1,3)$ entry is not $3/4$ or the QTL2 middle total is not
$-7/16$. \ev{our own calculation}

The earlier reading $B=-M$ with unsigned $G$ gives $1/4$ in that cell.
An intermediate reading, $B=+M$ with unsigned $G$, gives $7/8$. Neither is
$0.75$.

\subsection{The single-trait power on the earlier reading}

The $13\%$ figure is the single-trait detection rate for QTL3 on $Y_1$
under $B=-M$ and unsigned $G$, $100$ replicates, seed \texttt{20101027}.
It is recomputed by \texttt{simulation/gaps.py}, which writes
\texttt{simulation/out/gaps.json} and
\texttt{simulation/out/gaps\_table.csv}, committed in
\sha{c2769042e51806a77b2b70237fb4c014552e78fb}.
\begin{verbatim}
cd formalization/correlated-residuals/simulation
python3 gaps.py
\end{verbatim}
Under the glyph reading the same cell is $86.2\%$ at $1000$ replicates,
seed \texttt{20101041}, in that file. The printed value is $97$. The glyph
reading accounts for most of the difference between $13$ and $97$. It does
not meet $97$. \ev{our own simulation}

\subsection{Where the 2008 counts sit}

The 2008 conference paper, DOI \texttt{10.4148/2475-7772.1095}, SHA-256
\sha{ab2cc7ddc94025609cdaa83a6353b22dcdbb4693fedf2b70df71f5fc8aa1452a},
is a different design: $300$ lines and one QTL at $42$~cM. Its Table~2 text
layer gives multi-trait SEM overall power $91$, multi-trait least-squares
overall power $90$, and single-trait powers $2$, $39$, $4$, and overall
$42$. Trait-specific cells of the two multi-trait rows were not in the text
layer. Those counts are not a relabelling of SAGMB Table~1. The file is
recorded in \texttt{SOURCES.md}. Our three-QTL scan at $300$ lines, which is
not the conference design, gave SEM powers $89.3$, $85.0$, and $99.9$.
\ev{taken from a cited source, for the conference counts}

\section{What remains unresolved}
\label{sec:open}

The language in this section is graded on purpose. A difference between a
printed number and a number computed here is a source discrepancy under our
reading, or a gap this search did not close. It is not, by itself, evidence
that the printed number is a mistake.

\subsection{The gap between \(64/77/96\) and \(51/54/86\)}

Table~1 prints SEM powers $64$, $77$, and $96$ for the three QTLs.
The glyph baseline at the printed positions, $n=150$, threshold $14.86$,
and $1000$ replicates gives $50.8$, $53.5$, and $86.1$
(\texttt{simulation/out/assumptions/baseline.json}, seed
\texttt{20101101}, commit
\sha{82ff59d6bcfb02c273aaf1f83195e0d120fe0915}).
Rounded to the nearest percent, the comparison is $64/77/96$ against
$51/54/86$. The command is
\begin{verbatim}
cd formalization/correlated-residuals/simulation
python3 assumptions.py --configs baseline --reps 1000
\end{verbatim}
The absolute gaps are about $13$, $24$, and $10$ points. QTL2 is the large
one. A second seed stays in the same range. Runs that change one modeling
choice at a time are in Section~\ref{sec:gaps}. None of the choices
the simulation paragraph leaves open closes the QTL2 gap while keeping the
other two cells near the printed triple.
\ev{our own simulation, set next to a cited table}

Two calculations reproduce the middle cell and miss the first. At $n=250$
and the printed positions the SEM powers are $81.6$, $77.1$, and $98.7$.
The analytic approximation with the other QTL genotypes controlled, at
$n=150$, is about $81$, $77$, and $100$. Both put QTL1 near $80$ rather than
near $64$. Replacing $\Psi$ by $I$ gives about $84$, $80$, and $99$, which
is the same shape and contradicts equation~(8). An oracle scan that uses the
realized genotype gives about $65$, $64$, and $94$. A $\chi^2$ threshold on
$3$ degrees of freedom, which is not the printed $14.86$ on $4$ degrees of
freedom, gives $62.6$, $64.6$, and $92.9$ on the baseline.

\subsection{No text-consistent design matched all three numbers}

\texttt{qtl1\_search.py} varies the first QTL position, the sample size, a
marker gap around QTL1, local distortion, residual heterozygosity, a
two-plot line mean, and two cofactor rules that the simulation sentence does
not describe. Each specification is tagged with whether it is consistent
with the printed sentence. The search did not find a specification that is
consistent with that sentence and close to all three of $64$, $77$, and
$96$.

The smallest sum of squared errors among the percentage triples computed in
that search is the inconsistent point QTL1 at $24$~cM and $n=250$: SEM
powers $65.2$, $73.3$, and $99.3$, squared error about $26$, of which the
third cell is the larger share. The output is
\texttt{simulation/out/qtl1/qtl1\_24\_n250.json}, commit
\sha{b4e8a3c594f12f419560d5cb2a569b649ae93a55}.
\begin{verbatim}
cd formalization/correlated-residuals/simulation
python3 qtl1_search.py --configs qtl1_24_n250 --reps 1000
\end{verbatim}
Moving QTL1 off $22$~cM pulls the first power toward $64$ and pulls the
second down from $77$. The $99.3$ against a printed $96$, read as a
$\mathrm{Binomial}(100, 0.993)$ count, has lower-tail probability about
$0.006$ for a count of $96$ or fewer. That comparison is descriptive. It
does not identify a mechanism.

A global-maximum reading, one peak on the chromosome, produces powers on the
order of $16$, $14$, and $65$. The printed powers sum to $237$, so a rule
that returns only one detection cannot be the source of all three printed
numbers. Reading $14.86$ as a LOD threshold produces power $0$ at these
effect sizes.

\subsection{A probability on the order of \(10^{-12}\)}

Table~1 is a $100$-replicate experiment. Taking the baseline rates
$0.508$, $0.535$, and $0.861$ as the success probabilities,
\begin{quote}\noindent
\sha{python3 formalization/correlated-residuals/simulation/printed_probability.py}
\end{quote}
writes \texttt{simulation/out/printed\_probability.json}. The output is the
product of three binomial probability mass functions,
$\mathrm{Bin}(100, p_i)$ at $64$, $77$, and $96$. The product is
$1.37\times 10^{-12}$.

The product treats the three QTL counts as independent. They are computed
from the same replicates, so they are not independent. The figure is an
order of magnitude under that extra assumption. It says that a
$100$-replicate draw from the baseline rates is a poor match to the printed
triple. It does not say how the printed triple was obtained.
\ev{our own calculation, under an independence assumption we flag}

\subsection{The printed pair \(-0.435\) and \(0.065\)}
\label{sec:rounding}

Table~2, and Table~3 of the Genetics Research companion (the ARS file
Man851), print $-0.435$ and $0.065$ for the QTL2 effects that the glyph
arithmetic gives as $-7/16=-0.4375$ and $1/16=0.0625$. The same two
decimals appear in both papers we could read. Rounding $-0.4375$ to three
decimals is $-0.438$, not $-0.435$, and rounding $0.0625$ to three decimals
is $0.063$ or $0.062$ according to the tie rule, not $0.065$.

This is a source discrepancy under our reading of which cells those printed
numbers occupy. It is consistent with a rounding of a different intermediate
quantity, with a typesetting substitution, or with a coefficient we have not
reconstructed. The two papers agreeing with each other makes a one-time
typographical slip a weaker candidate, and it does not choose among the
others. \ev{source discrepancy (our reading)}

\subsection{SEM power above least-squares power}
\label{sec:semls}

On the recursive model the multi-trait SEM likelihood ratio and the
multi-trait least-squares likelihood ratio are the same statistic. In the
baseline profiles the maximum absolute difference of the mean profiles is
about $7\times 10^{-14}$. Table~1 nevertheless prints SEM powers $7$, $8$,
and $2$ points above the least-squares powers ($64$ against $57$, $77$
against $69$, $96$ against $94$). Figure~2 does not print a numeric scale
that would let those gaps be read off the plot.

Equal likelihood ratios make equal powers under a common threshold. A
difference of that size, repeated at all three QTLs, is a source discrepancy
under our reading of the two columns. It would also arise if the two columns
were computed under different thresholds or different cofactor sets. The
simulation paragraph describes one scan. The present runs do not separate
those possibilities. \ev{source discrepancy (our reading)}

\section{Code and source availability}
\label{sec:availability}

Each item was checked. The result is the result of that check.

\begin{itemize}
\item The 2008 conference paper (DOI \texttt{10.4148/2475-7772.1095})
contains no code and no code supplement. The software sentence names SAS or
Mplus. SHA-256 of the PDF we read:
\sha{ab2cc7ddc94025609cdaa83a6353b22dcdbb4693fedf2b70df71f5fc8aa1452a}.
\item The De Gruyter page for SAGMB article~38
(DOI \texttt{10.2202/1544-6115.1552}) showed no supplement, and the version
of record is paywalled. The file replicated here is the ARS manuscript
Man850, SHA-256
\sha{27c4f3926c9b5af9b399c9b3ef36918c28e7967ef777444b620279dd91628a81}.
The publisher file was not compared with that manuscript. The request for
the version of record is
\url{https://github.com/alectionary/vibe-math/issues/19}.
\item Mi's dissertation, \emph{Genetic applications using structural
equation modeling}, University of Nebraska--Lincoln, 2009, is restricted.
The catalog page is
\url{https://digitalcommons.unl.edu/dissertations/AAI3350451/}
and the ProQuest record is document \texttt{304940336}. The full text was
not obtained. The acquisition request is
\url{https://github.com/alectionary/vibe-math/issues/18}.
\item The Genetics Research companion (ARS file Man851, SHA-256
\sha{630a5e7b90665fd680d58915f6f9fe45e5051e8a5f6fff6af2cf16f8ae3b3739})
says that the programs were written in SAS PROC IML and are available by
email to \texttt{xjmixu@yahoo.com}. No message was sent. The request for
those programs, without sending the email, is
\url{https://github.com/alectionary/vibe-math/issues/20}.
\item Xiaojuan Mi works in industry as a contractor for Eli Lilly via
Tech Data Services, and was formerly at Duke Clinical Research Institute.
No academic webpage was found. Kent Eskridge's faculty page was not used as
a substitute for the code or for the dissertation.
\end{itemize}

\section{Next steps}
\label{sec:next}

Three steps are open, and this note does not carry them out.

The dissertation is the document most likely to contain the SAS code, the
marker list, and any appendix the articles omit. The route is an interlibrary
loan or another institutional copy. Filing that request is outside this
replication. The queue entry is issue~18, cited above.

A letter to Kent Eskridge, and a letter to Xiaojuan Mi, could ask for the
PROC IML programs and for the design details that decide among the readings
in Section~\ref{sec:open}. Those letters are to be drafted and sent by a
person. This replication does not send email and does not write the letters.

The formalization that this repository can continue without those documents
is local. The wheat fit is still blocked on the line identifiers, the ANOVA
residuals, the $\pm 1$ codes, and the edge set of Figure~5. The Drton,
Foygel, and Sullivant characterization of generically identifiable mixed
graphs, the half-trek gap of Foygel, Draisma, and Drton, and a polynomial-time
ancestral decomposition in the sense of Drton and Weihs are recorded as
separate questions (L-F145, L-F146, L-F147) and are not theorems here. A Lean
proof of the delta-method central limit theorem is the natural next
statement about the plug-in that is already proved consistent.

\section{The question, as they pose it}

The source is Mi et al., Statistical Applications in Genetics and Molecular
Biology~9 (2010), article~38. On printed page~18 (PDF page~20) they write:
\begin{quote}
Here, we assumed that the residual errors of the traits were independent
which we believe is a reasonable assumption since we removed the effects of
environment by using, for the $y$ values in equation~(2), the residuals from
a main effects linear model with environment and block within environment as
the main effects. However, even after removing environment main effects, the
traits may still be correlated and the trait error terms can be modeled as
being correlated with each other.
\end{quote}
\ev{taken from a cited source}
The parent request is that extension. It is not a numbered conjecture. Mi's
2009 dissertation (ProQuest AAI3350451) is the obvious place a fuller design
might have been written. It is restricted to a ProQuest preview. It was not
consulted.

\section{The restriction}

Equation~(7) is the only fully printed structural pattern. On PDF page~10
the equals sign ends at $x=145.9$ points and the parenthesis of $M$ begins at
$x=148.1$. No operator fits between them, so the display is
$y=My+Gx+\varepsilon$. The minus at $x=189.1$ is inside $M$, on the middle
row, and three further minuses sit on $G_{11}$, $G_{22}$, and $G_{32}$.
\ev{our reading of the glyph positions}
\[
M=\begin{pmatrix}0&1/2&1/4\\0&0&-1/2\\0&0&0\end{pmatrix},\qquad
G=\begin{pmatrix}-1/8&1/4&1/2\\1/2&-1/2&1/2\\1/4&-1/8&1/4\end{pmatrix}.
\]
Equation~(8) is $\Psi=\mathrm{diag}(1.6,1.8,2.5)$.
\ev{taken from a cited source}
In $y=By+Gx+\varepsilon$ one then has $B=M$. A first reading placed the
interior minus in front of $M$ and took every entry of $G$ as positive.
That reading is the population of the first replication below. The zero
pattern is the same either way: the free positions are still $(1,2)$,
$(1,3)$, and $(2,3)$.
The wheat Figure~5 graph has four traits and is not this matrix. The
restriction is covariance identification for this three-trait recursive
pattern: diagonal $\Psi$, one extra residual covariance, and the three ways
to drop one edge and free the residual pair that edge had used. Singular
$\Psi$ is outside the claim. The general mixed-graph characterization is not
this restriction.

\section{What we proved}

The following are Lean proofs on the positive-definite domain. The axioms are
\texttt{propext}, \texttt{Classical.choice}, and \texttt{Quot.sound}. There is
no \texttt{sorry}. \ev{our own Lean proof}
\begin{itemize}
\item \texttt{recover\_sigma}: diagonal residuals and positive-definite $\Phi$
are a left inverse of the joint covariance on the equation-(7) support.
\item \texttt{bow12\_same\_law} and its two siblings: each one-edge residual
extension collides, for every exogenous $\Gamma$.
\item \texttt{recover\_bow}, \texttt{recover\_bow01}, \texttt{recover\_bow23}:
each bow-free missing-edge pattern is a left inverse.
\item \texttt{unitUpper\_factor\_unique} and \texttt{sigmaRec\_injective}:
in every trait and covariate dimension, a recursive pattern with positive
diagonal residuals has a unique unitriangular factorization, and the joint
blocks determine $(B,\Psi,\Gamma,\Phi)$.
\item \texttt{plug\_in\_consistent\_ae} and the three bow-free analogues
\texttt{plug\_in\_bow\_consistent\_ae},
\texttt{plug\_in\_bow01\_consistent\_ae},
\texttt{plug\_in\_bow23\_consistent\_ae}: pairwise i.i.d.\ samples with
integrable coordinate products make those plug-ins almost surely consistent.
\end{itemize}
The Gaussian delta-method formula for the diagonal plug-in is a written
derivative together with Isserlis' covariance, checked numerically in
\texttt{delta\_method.py}. \ev{our own written proof, and our own simulation}
It is not a Lean central limit theorem, and it is not listed among the proved results.

\begin{center}
\footnotesize
\begin{tabular}{@{}p{0.24\textwidth}p{0.34\textwidth}p{0.36\textwidth}@{}}
\toprule
Statement & Lean declaration & File and lines \\
\midrule
Diagonal left inverse & \texttt{recover\_sigma} & \href{\proofbase/CorrelatedResiduals/Diagonal.lean\#L246-L279}{\texttt{Diagonal.lean}, 246--279} \\
One-edge collision & \texttt{bow12\_same\_law} and two siblings & \href{\proofbase/CorrelatedResiduals/Collision.lean\#L114-L160}{\texttt{Collision.lean}, 114--160} \\
Bow-free left inverse & \texttt{recover\_bow} & \href{\proofbase/CorrelatedResiduals/BowFree.lean\#L187-L213}{\texttt{BowFree.lean}, 187--213} \\
Other missing edges & \texttt{recover\_bow01}, \texttt{recover\_bow23} & \href{\proofbase/CorrelatedResiduals/BowMissing.lean\#L148}{\texttt{BowMissing.lean}, 148 and 353} \\
Triangular factor & \texttt{unitUpper\_factor\_unique} & \href{\proofbase/CorrelatedResiduals/Triangular.lean\#L68-L135}{\texttt{Triangular.lean}, 68--135} \\
Recursive identification & \texttt{sigmaRec\_injective} & \href{\proofbase/CorrelatedResiduals/Triangular.lean\#L298-L317}{\texttt{Triangular.lean}, 298--317} \\
Almost-sure consistency & \texttt{plug\_in\_consistent\_ae} and \texttt{plug\_in\_bow\_consistent\_ae} & \href{\proofbase/CorrelatedResiduals/Law.lean\#L174-L199}{\texttt{Law.lean}, 174--199}; \href{\proofbase/CorrelatedResiduals/BowConsistency.lean\#L270-L304}{\texttt{BowConsistency.lean}, 270--304} \\
\bottomrule
\end{tabular}
\end{center}

The links are pinned to commit \sha{05d69ea141adf3d5ee936d3b48061bde241250dc}, which contains these Lean files. The simulation script and the JSON outputs are commit \sha{5cd3c9f6d2f6e38f25e32146d74d55f3778436e1}. Appendix~\ref{app:lean} prints the central proof bodies.

Brito and Pearl (Structural Equation Modeling, 2002) prove that a recursive
linear model is identifiable when no residual covariance joins a variable to
its direct effect. Our three bow-free left inverses are the three-trait case
of that pattern, proved from the printed matrices rather than cited as a
substitute for the algebra. Drton, Foygel, and Sullivant (Annals of
Statistics, 2011) give a global criterion for mixed graphs; the half-trek
criterion of Foygel, Draisma, and Drton (Annals of Statistics, 2012) is a
generic sufficient condition. Neither is re-proved here, and the 2011
characterization stays a note with no inventory id. Li et al.\ (PLoS
Genetics, 2006) and Valente, Rosa, de los Campos, Gianola, and Silva
(Genetics, 2010) are earlier structural and recursive treatments of multiple
traits. Hidalgo-Contreras, Salinas-Ruiz, Eskridge, and Baenziger (Agronomy,
2021) fit a later Nebraska wheat SEM; their lines are a 2013 F6 panel, not
the 1999--2001 RICLs, and none of their estimates enter the simulations.

\section{Replication design}

The printed design, PDF pages~9--10, is 100 replicates of 150 RIL lines, one
chromosome of 100~cM, 11 evenly spaced markers, additive QTLs at 22, 42, and
78~cM, phenotypes from equations~(7)--(8), and a test every 1~cM. QTL
genotypes are coded $+1$ and $-1$. Cofactors are all markers except the
flanking markers. The printed critical values are 14.86 and 10.60.
\ev{taken from a cited source}
Those two numbers are the chi-square $0.005$ quantiles on 4 and 2 degrees of
freedom ($14.860$ and $10.597$). The paragraph on PDF page~9 writes the
Bonferroni approximation $\chi^2_{\alpha/M,\,m+1}$. Ten marker intervals and
$\alpha=0.05$ give the tail probability $0.005$, with $m+1=4$ for three
traits and $m+1=2$ for one trait. \ev{our own calculation, from the cited formula}

The paper does not state the following, so each is our assumption.
\ev{our assumption}
\begin{enumerate}
\item The map function is Haldane, with no interference.
\item An RIL is an F1 selfed for 12 generations. A leftover heterozygote is
coded $\pm 1$ with equal probability. There is no segregation distortion.
\item The eleven markers are at $0,10,\ldots,100$~cM.
\item The analysis uses the centered covariance with divisor $n$, which is
the fitting function~(6). We did not run SAS PROC TCALIS.
\item The flanking expectation is the conditional mean of the QTL genotype
under that same 12-generation law. The 10~cM window mentioned later in the
paper is for the wheat scan, not for this simulation, and was not used.
\item A named QTL is detected when the maximum likelihood ratio in
$[0,32]$, $(32,60]$, or $(60,100]$ exceeds the threshold. Effects are read
at that argmax.
\end{enumerate}
Seeds are \texttt{20101027} for this study. The outputs are
\texttt{study1.json}, \texttt{study1\_power.csv}, and
\texttt{study1\_effects.csv} under \texttt{simulation/out}.

For a complete recursive order, the Gaussian likelihood of the diagonal
structural model equals the likelihood of the multivariate regression on the
same exogenous regressors: the sum of the successive residual log variances
is the log determinant of the residual covariance. The two joint likelihood
ratios are therefore the same number. In the run, the maximum absolute
difference of the two mean profiles is $7\times 10^{-14}$.
\ev{our own written proof, confirmed by our own simulation}
Mi et al.\ describe the two profiles as very close (PDF page~14). Their
Table~1 powers still differ by 7, 8, and 2 replicates out of 100.

\section{Replication results}

Table~\ref{tab:power} compares detection percentages. Our multi-trait powers
keep the printed order, third QTL above the first two, and they sit 5 to 16
points below Table~1. The single-trait column for the first trait at the
first QTL matches Table~1 exactly (6\%). The single-trait column for the
first trait at the third QTL does not: Table~1 prints 97\% and the
replication gives 13\%.

\begin{table}[ht]
\centering
\caption{Detection percentages. ``Mi'' is Table~1 on PDF page~12. ``Ours''
is 100 replicates, seed \texttt{20101027}. SEM and LS coincide in our run.}
\label{tab:power}
\begin{tabular}{crrrrrr}
\hline
QTL & SEM Mi & SEM ours & LS Mi & ST $Y_1$ Mi & ST $Y_1$ ours & ST overall ours \\
\hline
1 & 64 & 59 & 57 & 6 & 6 & 24 \\
2 & 77 & 61 & 69 & 3 & 0 & 35 \\
3 & 96 & 86 & 94 & 97 & 13 & 34 \\
\hline
\end{tabular}
\end{table}

The algebraic total effects under $B=-M$ are, by columns of $G$,
\[
(I+M)^{-1}G
=\begin{pmatrix}-1/8&0&1/4\\ 3/8&7/16&3/8\\ 1/4&1/8&1/4\end{pmatrix}.
\]
\ev{our own calculation}
Table~2's parameter block matches this matrix on $Y_2$ and $Y_3$ for the
first and third QTLs, and it does not match it on $Y_1$. The printed total
effect of the third QTL on $Y_1$ is $0.75$, against $0.25$ from this first
reading. Section~\ref{sec:gaps} tests that disagreement directly.

Direct-effect means at the SEM argmax are close to $G$ where the signal is
large. For the third QTL they are $0.460$ ($0.175$), $0.476$ ($0.179$), and
$0.242$ ($0.190$), against the printed SEM estimates $0.479$ ($0.176$),
$0.527$ ($0.166$), and $0.257$ ($0.205$) and against $G=(1/2,1/2,1/4)$.
\ev{our own simulation, set next to a cited source}
Position means are $22.65$ ($7.28$), $41.26$ ($5.25$), and $79.85$ ($3.76$),
against the printed $21.57$ ($2.83$), $41.95$ ($2.58$), and $78.03$ ($2.16$).
The biases are small. Our standard deviations are larger. The mean profile
exceeds $14.86$ only at 40~cM and on 75--82~cM. Figure~2 is described as one
simulated population, with excursions on 16--28, 38--48, and 70--88~cM. An
average profile is a different object, and ours is lower.

\section{Which hypothesis closes which gap}
\label{sec:gaps}

Five changes were rerun from fixed seeds. The full grid is
\texttt{simulation/out/gaps.json} and \texttt{gaps\_table.csv}. Closure is
$1-|\hat p-p_{\mathrm{printed}}|/|p_{\mathrm{first}}-p_{\mathrm{printed}}|$,
using the first reading's $100$-replicate powers as the denominator. A value
of $1$ means the printed number is met. A negative value means the estimate
moved away from the table. \ev{our own simulation}

\begin{table}[ht]
\centering
\footnotesize
\caption{Gap closure. SEM powers are the three QTLs. ST$_3$ is single-trait
power for QTL3 on $Y_1$ (printed 97). ST$_2$ is single-trait power for QTL2
on $Y_1$ (printed 3). SDs are the three SEM position standard deviations.
$T_{31}$ is the algebraic total effect of QTL3 on $Y_1$ (printed 0.75).}
\label{tab:gaps}
\begin{tabular}{lrrrrrrr}
\hline
Reading & SEM & ST$_3$ & ST$_2$ & position SD & $T_{31}$ & ST$_3$ closed & $T_{31}$ closed \\
\hline
Printed & 64, 77, 96 & 97 & 3 & 2.83, 2.58, 2.16 & 0.75 & --- & --- \\
First, $n=100$ & 59, 61, 86 & 13 & 0 & 7.28, 5.25, 3.76 & 0.25 & 0 & 0 \\
First, $n=1000$ & 65, 63, 87 & 9 & 1 & 6.20, 5.32, 4.38 & 0.25 & $-0.05$ & 0 \\
$B=+M$, $G$ unsigned & 68, 57, 86 & 87 & 44 & 5.90, 5.33, 4.60 & 0.875 & 0.88 & 0.75 \\
Glyph reading, $n=100$ & 46, 45, 80 & 85 & 1 & 5.57, 5.93, 5.10 & 0.75 & 0.86 & 1 \\
Glyph, $n=1000$ & 52, 48, 87 & 86 & 1 & 5.15, 5.42, 3.47 & 0.75 & 0.87 & 1 \\
Glyph, no cofactors & 46, 65, 100 & 100 & 65 & 6.99, 9.12, 3.80 & 0.75 & 0.96 & 1 \\
Glyph, wheat cofactors & 34, 24, 89 & 95 & 1 & 7.33, 7.91, 3.77 & 0.75 & 0.98 & 1 \\
Glyph, $\Psi=I$ & 80, 78, 98 & 97 & 2 & 4.86, 4.02, 2.79 & 0.75 & 1.00 & 1 \\
Glyph, window $\pm 5$ cM & 43, 41, 80 & 83 & 1 & 2.80, 2.76, 2.03 & 0.75 & 0.83 & 1 \\
Glyph, $\alpha=0.05$ & 77, 84, 98 & 95 & 14 & 5.57, 5.93, 5.10 & 0.75 & 0.98 & 1 \\
\hline
\end{tabular}
\end{table}

\textbf{Sign convention.} $B=+M$ with unsigned $G$ gives total effect $0.875$
on $Y_1$ for QTL3, not $0.75$. The glyph reading gives
\[
(I-M)^{-1}G
=\begin{pmatrix}1/8&0&3/4\\ 3/8&-7/16&3/8\\ 1/4&-1/8&1/4\end{pmatrix}.
\]
\ev{our own calculation}
Every Table~2 parameter cell matches this arithmetic except the QTL2 row for
$Y_2$, where the exact values are $-7/16=-0.4375$ and $1/16=0.0625$ and the
table prints $-0.435$ and $0.065$.
\ev{source discrepancy (our reading)}
The first reading's total of $0.25$ is what produced the single-trait power
$13\%$. Under the glyph reading, $1000$ replicates (seed \texttt{20101041})
give $86.2\%$ (Clopper--Pearson interval $83.9$--$88.3$). That closes most of
the $97$ versus $13$ gap and does not meet $97$.

\textbf{Scan method.} The simulation paragraph on PDF page~10 fits every
marker except the flanking pair, and the methods section replaces the QTL
genotype by its flanking-marker expectation (Haley and Knott, 1992).
Dropping the cofactors, on the glyph reading, raises the single-trait QTL3
power to $100\%$ and the SEM power for QTL3 to $100\%$, and it leaves the
first SEM power at $46$. The wheat rule on PDF page~11, a $P<0.1$ screen
with at most five cofactors and a $10$~cM window, gives single-trait power
$95\%$ (interval $88.7$--$98.4$, which contains $97$) and SEM powers $34$,
$24$, and $89$. Nearest-marker regression gives the same pattern as dropping
cofactors. None of these is the sentence printed for the simulation, and
none of them brings all three SEM powers onto $64$, $77$, and $96$.

\textbf{Threshold and window.} One hundred phenotype shuffles of a glyph-reading
replicate have $95\%$ quantile $14.74$ for the maximum likelihood ratio,
next to the printed $14.86$. A permutation threshold at $\alpha=0.05$ does
not move the powers. The $500$ permutations in the paper are in the wheat
section. Replacing $14.86$ and $10.60$ by the $\alpha=0.05$ quantiles on 4
and 2 degrees of freedom ($9.49$ and $5.99$) raises the glyph-reading SEM
powers to $77$, $84$, and $98$ and the single-trait QTL3 power to $95$. That
threshold is not the printed one. Taking the peak within $5$~cM of the true
site, under the glyph reading, gives position standard deviations $2.80$,
$2.76$, and $2.03$, against $2.83$, $2.58$, and $2.16$. The midpoint windows
used in the first replication were our assumption. The paper does not name
the window that assigns a peak to a named QTL.

\textbf{Unprinted design.} Eleven evenly spaced markers on a segment of
$100$~cM, with markers at the ends, are at $0,10,\ldots,100$. Shifting them
to $5,15,\ldots,95$ does not close the SEM gap. Standardizing each trait to
sample standard deviation $1$ leaves the likelihood ratio essentially
unchanged. Replacing $\Psi$ by $I$ does close the single-trait QTL3 gap
($97\%$) and puts the second and third SEM powers on $78$ and $98$, while the
first SEM power rises to $80$ (interval $70.8$--$87.3$, which does not
contain $64$). The simulation sentence says the residuals are drawn from
matrix~(8). $\Psi=I$ contradicts that sentence. It is a sensitivity, not the
display.

\textbf{Monte Carlo error.} At $1000$ replicates the first reading gives SEM
powers $65.1$, $62.5$, and $86.6$, with intervals $62.1$--$68.1$,
$59.4$--$65.5$, and $84.3$--$88.7$. The printed $64$ lies in the first
interval. The printed $77$ and $96$ do not. Single-trait power for QTL2 on
$Y_1$ is $1.2\%$ (interval $0.6$--$2.1$) under that reading and $0.8\%$
($0.3$--$1.6$) under the glyph reading. The printed $3$ is inside the
$n=100$ interval for zero detections ($0$--$3.6$) and outside the
$n=1000$ intervals. The discrepancy is three replicates in Table~1.

The most likely account is the glyph reading of equation~(7), together with
a peak window of about $5$~cM that the paper does not state. That reading
removes the $0.75$ discrepancy and most of the single-trait QTL3 discrepancy,
and the narrow window removes the position-SD discrepancy. It does not
remove the SEM gap for QTL2 and QTL3. Those two percentages move onto the
printed values only in runs that drop the printed cofactors or replace
matrix~(8). Both of those runs conflict with a sentence on PDF page~10.

\section{Explanations considered}
\label{sec:explanations}

The gaps that remained after the glyph reading, on a fresh 1000-replicate
run (seed \texttt{20101101}), are SEM powers $50.8$, $53.5$ and $86.1$
against $64$, $77$ and $96$, single-trait power for QTL3 on $Y_1$ of $88.6$
against $97$, and single-trait power for QTL2 on $Y_1$ of $1.2$ against $3$.
The intervals are $50.4$--$56.6$, $83.8$--$88.2$, $86.5$--$90.5$ and
$0.6$--$2.1$. A second seed (\texttt{20101102}) gives SEM powers $49.4$,
$48.3$ and $84.6$. The QTL2 gap of about $24$ points is not a seed effect.
Family A changes one modeling choice at a time. Family B compares the
printed counts with the arithmetic, with a noncentral chi-square
approximation, and with the companion designs. The grid is
\texttt{simulation/out/assumptions/}. \ev{our own simulation}

\begin{table}[ht]
\centering
\scriptsize
\caption{Explanations considered. Closure is the fraction of the
seed-\texttt{20101101} gap removed. The four fractions are QTL2 SEM,
QTL3 SEM, single-trait QTL3 on $Y_1$, and single-trait QTL2 on $Y_1$.}
\label{tab:explanations}
\begin{tabular}{p{3.3cm}cp{3.6cm}p{2.4cm}p{2.6cm}}
\hline
Hypothesis & Family & Evidence & Gap closed & Verdict \\
\hline
Infinite RIL, $2r/(1+2r)$ & A & SEM $56,55,91$ & $.07/.54/.36/-.06$ & open, and too small \\
Kosambi, infinite RIL & A & SEM $56,54,91$ & $.03/.49/.31/-.06$ & open in SAGMB; not Table 1 \\
Flanking formula & A & within the seed range & about $0$ & cannot be distinguished from the 12-generation expectation \\
Per-equation tests & A & any equation above $10.60$ gives $52,52,78$ & negative on QTL3 & does not match \\
Permutation threshold & A & $95\%$ quantile $16.1$ in 200 shuffles & power falls & $14.86$ is slightly liberal; the gap is not a low threshold \\
$\Psi=I$ & A & SEM $84,80,99$; ST$_3=98$ & $.89/.70/.87/.06$ & contradicts matrix (8) \\
Entries used as SDs & A & power falls to $24,25,59$ & negative & not the covariance sentence \\
Half contrast, $0/1$ & A & power falls to $11,9,23$ & negative & the sentence says $+1$ and $-1$ \\
F2 or backcross & A & power falls & negative & the sentence says RIL \\
Oracle genotype & A & SEM $65,64,94$; ST$_3=95$ & $.46/.85/.80/.94$ & not the flanking expectation \\
QTLs on markers & A & SEM $72,73,96$ & $.84/.97/.85/.22$ & contradicts $22,42,78$ \\
No cofactors & A & ST$_2$ rises to $66$ & $.59/.65/.70$ and ST$_2$ breaks & contradicts the cofactor sentence \\
$n=250$ & B & SEM $82,77,99$; $Y_2$ powers $31,38,32$ & $1.00/.73/.76/.56$ & matches QTL2, not QTL1 \\
$n=300$ & B & SEM $89,85,100$ & overshoots QTL2 & above Table 1 \\
Chi-square, QTLs controlled & B & powers $81,77,100$ at $n=150$ & QTL2 matches & not the marker-cofactor scan \\
\hline
\end{tabular}
\end{table}

No choice the simulation paragraph leaves open closes the QTL2 gap. The
infinite-RIL formula is the one printed in the 2008 conference note, and
Kosambi is unnamed in SAGMB; together they move QTL2 from $53.5$ to $54.2$.
The no-double-crossover expectation in that note, and an F2 expectation
applied to RIL data, stay inside the seed range of the baseline. A
$95\%$ permutation threshold is $16.1$, and the median of $40$ per-replicate
thresholds, each the maximum of $19$ null scans, is $16.9$. Both are above
$14.86$, so a permutation rule makes the powers smaller ($43.5$, $44.2$,
$82.0$ at $16.1$). Replacing $14.86$ by the $\alpha=0.05$ quantile $9.49$
raises the baseline to $84.7$, $82.6$ and $98.5$, which is not the printed
critical value. Standardizing the traits, scanning every $2$ or $5$~cM, and
shifting the eleven markers off the ends do not close QTL2. The $\pm 1$ and
$0/2$ codes have the same homozygote contrast, so they have the same
likelihood ratio; a $0/1$ code with the printed $G$ halves that contrast and
the power collapses. \ev{our own simulation}

The printed $77$ is consistent with two calculations that are not the
printed marker scan. With the three QTL genotypes in each structural
equation, the noncentral chi-square approximation at $n=150$ and threshold
$14.86$ gives joint powers $81\%$, $77\%$ and $100\%$. The partial
single-trait power for QTL3 on $Y_1$ is $99.7\%$, and for QTL2 on $Y_1$ it
is $0.1\%$. The same approximation at $n=250$ and $n=300$ is above $97\%$
for every QTL. \ev{our own calculation}
The marker-cofactor scan at $250$ lines, the sample size of Mi et al.\
Genetics Research (2010), gives SEM powers $81.6$, $77.1$ and $98.7$,
single-trait $Y_1$ powers $3.2$, $2.2$ and $99.0$, and single-trait $Y_2$
powers $31$, $38$ and $32$ against the printed $29$, $43$ and $33$. QTL2's
SEM cell matches. QTL1's SEM cell is $82$ against $64$. At $300$ lines, the
sample size of the 2008 conference note, the three-QTL scan is $89.3$,
$85.0$ and $99.9$, above Table 1. \ev{our own simulation}

Those companion designs are not Table 1. Genetics Research simulates $250$
lines and two QTLs, at $42$ and $78$~cM, with codes $2$ and $0$, and the
detection rule is a posterior peak rather than the likelihood-ratio
threshold. Its Table 2 powers are $99$ and $100$. Equation~(15) and Figure~2
of that paper print the same path signs as the glyph reading, $1/2$, $1/4$
and $-1/2$, and the same rounded QTL2 row $-0.435$ and $0.065$.
\ev{taken from a cited source}
The conference note simulates $300$ lines and one QTL at $42$~cM, names
Haldane's map and the expansion $2r/(1+2r)$, and uses codes $2$ and $0$.
Its Table~2 text layer gives multi-trait SEM overall power $91$, multi-trait
least-squares overall power $90$, and single-trait powers $2$, $39$, $4$,
and overall $42$. Trait-specific cells of the two multi-trait rows were not
in the text layer. Neither table is a relabeling of SAGMB Table~1.
\ev{taken from a cited source}

Table 2's QTL2 row prints $-0.435$ and $0.065$ where the path product is
$-7/16=-0.4375$ and $1/16=0.0625$. The same printed decimals appear in the
Genetics Research table. \ev{source discrepancy (our reading)}
On a complete recursive model the multi-trait SEM and multi-trait LS
likelihood ratios are the same number, and in every run here their powers
agree. Table 1 prints SEM powers $7$, $8$ and $2$ points above the LS
powers, while the text says the two profiles are very close.
\ev{source discrepancy (our reading)}
Figure 2 is described as one simulated population and does not print a
numeric likelihood-ratio scale, so the power counts cannot be read off the
figure.

The best account is that no open choice in the simulation paragraph
produces Table 1, and that the cell $77$ is consistent both with the
chi-square power of these direct effects at $150$ lines when the other QTL
genotypes are controlled and with the marker scan at $250$ lines. Both of
those readings put the first QTL near $80\%$ rather than at the printed
$64\%$. The single-trait QTL2 cell, $3$ against our $1.2$, is consistent
with the algebraic total effect of zero.

\subsection{A joint reading of 64, 77 and 96}
\label{sec:qtl1}

A second block asked which reading could produce the SEM triple $64$, $77$,
$96$ together. Every entry of equation~(7) was measured again on PDF page 10.
The minus signs are four glyphs. One ends at $x=250.7$ on $G_{11}=-1/8$.
One ends at $x=194.5$ on $M_{23}=-1/2$. Two end at $x=316.7$, on
$G_{22}=-1/2$ and $G_{32}=-1/8$. No other entry carries a minus, and there is
no minus between the equals sign and $M$. The Table~1 digits are the
characters 6 and 4, then 77, then 96. \ev{taken from a cited source}
The scans below use $1000$ replicates. Squared error is the sum of squared
differences of the three SEM percentages. The chi-square treats each printed
percentage as a count out of $100$ draws from the simulated rate.
\ev{our own simulation}

\begin{table}[ht]
\centering
\scriptsize
\caption{SEM triples ranked by squared error against $64$, $77$, $96$.
The chi-square is for those three cells. A text-consistent row is one the
simulation paragraph on PDF page 10 states.}
\label{tab:qtl1}
\begin{tabular}{p{4.4cm}crrr p{3.2cm}}
\hline
Hypothesis & Family & SEM powers & SSE & $\chi^2$ & Consistent with the text \\
\hline
QTL1 at 24~cM, 250 lines & A+C & 65, 73, 99 & 26 & 16.4 & no \\
QTL1 at 25~cM, 250 lines & A+C & 62, 69, 99 & 74 & 13.8 & no \\
QTL1 at 23~cM, 250 lines & A+C & 72, 73, 99 & 80 & 8.6 & no \\
Oracle genotype, 150 lines & A & 65, 64, 94 & 163 & 7.4 & no \\
Threshold $\chi^2_{3,0.995}$ & B & 63, 65, 93 & 165 & 8.3 & no; printed value is 14.86 \\
Marker gap around QTL1, 250 lines & A+C & 78, 76, 98 & 212 & 13.6 & no \\
Printed positions, 250 lines & C & 82, 77, 99 & 317 & 26.3 & no; the sentence says 150 \\
Two plots averaged, $\Psi$ at plot level & A & 63, 59, 100 & 341 & 174 & no \\
Local allele frequency 0.30 at QTL1 & A & 62, 61, 85 & 392 & 20.8 & no \\
$\Psi=I$ & A & 84, 80, 99 & 416 & 39.3 & no; contradicts (8) \\
Glyph baseline, 150 lines & A & 51, 54, 86 & 824 & 37.4 & yes \\
One peak on the chromosome & B & 16, 14, 65 & 4602 & --- & the printed sum is 237 \\
14.86 applied to the LOD & B & 0, 0, 0 & --- & --- & 14.86 is the chi-square value \\
$G$ transposed & A & 55, 94, 19 & 6308 & 441 & no; columns are labeled QTL \\
\hline
\end{tabular}
\end{table}

The smallest squared error is $65.2$, $73.3$, $99.3$. Single-trait powers on
$Y_1$ are $2.6$, $1.9$, $99.1$ against $6$, $3$, $97$, and on $Y_2$ they are
$22$, $37$, $33$ against $29$, $43$, $33$. QTL1's $65$ cannot be
distinguished from the printed $64$ at the precision of $100$ replicates.
QTL2 remains about four points low. QTL3 at $99.3$ makes a printed $96$
unlikely if that rate were true: the lower-tail probability is about
$0.006$. The position $24$~cM and the sample size $250$ are not the
simulation sentence. Moving QTL1 from $22$ to $24$~cM is what pulls the
first power down from the $82$ of the $250$-line scan at the printed
positions. The same move pulls the second power down from $77$ to $73$.
The two cells do not land on $64$ and $77$ together.
\ev{our own simulation}

Family A does not supply a text-consistent joint match. Flipping the sign
of $G_{11}$, which the remeasured glyph does not support, gives $55$, $50$,
$85$. Halving $G_{21}$ gives $17$, $53$, $87$. A residual heterozygosity of
$0.15$ gives $20$, $18$, $57$. Placing QTL1 $10$~cM from QTL2 gives $24$,
$22$, $84$. Kosambi interference was already inside a few points of the
baseline. Standardizing the traits was $51$, $48$, $87$.

Family B does not supply the triple either. Stepwise cofactors and a
two-stage refit lower the first two powers and raise the third, to about
$33$, $28$, $98$ and $29$, $28$, $98$. A marker-bracket window, a $20$~cM
separation rule, and dropping the terminal $10$~cM leave the baseline near
$51$, $54$, $86$. The $\chi^2_3$ threshold $12.84$ gives $63$, $65$, $93$:
the first cell matches, and the second does not. The printed critical value
is $14.86$, the $\chi^2_4$ quantile. Reading $14.86$ as a LOD multiplies the
likelihood-ratio threshold by $2\log 10$ and the powers fall to zero. One
peak per chromosome assigns the global maximum to a single QTL, so the three
powers sum to at most $100$. On the baseline they are $16$, $14$, $65$.
The printed powers sum to $237$, so that counting rule cannot be the source
of all three cells. The flanking regressor has positive variance at every
interior test point on this marker grid, so a singular fit does not arise
and counting non-convergence as a miss does not change the powers.

Family C. The probability of the printed counts $64$, $77$, $96$ as three
independent $\mathrm{Bin}(100)$ draws from the baseline rates $0.51$,
$0.54$, $0.86$ is about $10^{-12}$. Swapping the SEM and LS columns of
Table~1 produces $57$, $69$, $94$, which is not a row of any scan here.
The PDF prints the digit 6, not 8, so $64$ is not an OCR reading of $84$.
The $250$-line scan at the printed positions is $82$, $77$, $99$. That
matches the cell $77$. It is near $84$, $77$, $96$ in the first cell and
not near $64$. The Genetics Research powers are $99$ and $100$ under a
different design. The $300$-line three-QTL scan is $89$, $85$, $100$.

The best joint account is negative. No tested reading that the simulation
paragraph states produces $64$, $77$, and $96$ together. The closest
percentage triple uses a QTL position and a sample size that paragraph does
not state, and its third cell is still high relative to $96$. What remains
is that gap: the first cell can be brought near $64$ by weakening QTL1 or
by a lower threshold, and the second cell can be brought to $77$ by using
$250$ lines or by controlling the other QTL genotypes, and those two
adjustments have not been found together in one design the text states.

\section{Correlated residuals}

The bow drops the edge into trait~0 from trait~2, sets that structural
coefficient to zero, and frees the residual covariance $c=\rho\sqrt{1.6\cdot 2.5}$.
The other equation-(7) coefficients stay. The exogenous variables are the
three QTL genotypes on the RIL map, not a marker scan.
\ev{our assumption, matching \texttt{recoverBow}}
The sample is the centered second-moment matrix, $n=150$, $400$ replicates,
seed \texttt{20101028}.

At the population covariance the bow plug-in recovers the parameter to
machine precision. The diagonal plug-in does not. Its probability limit for
the missing edge is $0.8\rho$: zero at $\rho=0$, $0.20$ at $\rho=0.25$, and
$0.48$ at $\rho=0.6$. The QTL coefficient $G_{11}=1/8$ has probability limit
$0.075$ at $\rho=0.25$ and $0.005$ at $\rho=0.6$. The coefficient $G_{13}=1/2$
moves to $0.45$ and $0.38$. \ev{our own simulation}
A Wald test that the missing edge is zero, using the diagonal delta-method
standard error, rejects at rate $0.058$ when $\rho=0$ and at rate $0.728$,
$0.877$, and $1$ at $\rho=0.25$, $0.30$, and $0.60$. The joint nine-degree
Wald test of all QTL coefficients is only mildly inflated at $\rho=0.25$
(rejection rate $0.065$) and reaches $0.297$ at $\rho=0.6$. The damage shows
up first as bias in two coefficients and as a spurious edge, not as an
immediate explosion of the joint QTL test.

\texttt{recoverBow} is centered on the QTL and path coefficients. At
$\rho=0.25$ the Monte Carlo means of $b_{12}$, $b_{23}$, $G_{11}$, and
$G_{13}$ are $-0.507$, $-0.498$, $0.127$, and $0.509$. Coverage of the
nominal 95\% delta-method intervals, with the Gaussian Isserlis covariance
plugged into the sample, is between $0.915$ and $0.958$ across the reported
coefficients and the six correlations. The genotypes are $\pm 1$, so
Isserlis is an approximation. \ev{our assumption inside the standard error}
The coverage near 95\% says the approximation is tolerable at $n=150$ for
these coefficients. Variance-component means sit a little low ($1.54$ against
$1.6$ at $\rho=0$) because the likelihood uses divisor $n$; the probability
limit is still the population value.

The same seeds were rerun with the glyph-reading coefficients: $b_{12}=+1/2$,
$b_{23}=-1/2$, the edge $b_{13}$ removed from $+1/4$, and the signed $G$.
\ev{our own simulation}
At $\rho=0.25$ the diagonal probability limit of the missing edge is again
$0.20$ and the Wald rejection rate is $0.728$. \texttt{recoverBow} means are
$b_{12}=0.493$, $b_{23}=-0.498$, $G_{11}=-0.123$, and $G_{22}=-0.508$, with
coverage between $0.920$ and $0.955$. The sign reading changes the parameter
values and does not change the bias or the coverage.

\section{Calibration to three phenotypic correlations}

Campbell et al., Crop Science~43 (2003), Table~2, print phenotypic
correlations for 95 recombinant inbred chromosome lines pooled across seven
environments. There are no standard deviations and no environment or block
adjustment. \ev{taken from a cited source, with the limits just stated}
The text layer of the copy we could read dropped minus signs. The paragraph
above the table says that SPSM and kernels per square meter were negatively
correlated with thousand-kernel weight. We therefore use
$\mathrm{Corr}(\mathrm{TKWT},\mathrm{SPSM})=-0.44$. The cells
$\mathrm{Corr}(\mathrm{TKWT},\mathrm{KPS})=0.08$ and
$\mathrm{Corr}(\mathrm{KPS},\mathrm{SPSM})=0.59$ are the signs in that text
layer; the prose does not restate them. \ev{our assumption for those two signs}
GYLD is the fourth trait in the application and is not in equation~(7). The
three-trait order is $Y_1=\mathrm{TKWT}$, $Y_2=\mathrm{KPS}$,
$Y_3=\mathrm{SPSM}$, which is the chain used in the indirect-effect
arithmetic on PDF page~17 of Mi et al. \ev{our assumption}

Treating the correlation matrix as a covariance, the diagonal recursive
factorization is
\[
B=\begin{pmatrix}0&0.521&-0.747\\0&0&0.590\\0&0&0\end{pmatrix},\qquad
\Psi=\mathrm{diag}(0.629,0.652,1).
\]
\ev{our own calculation}
The eigenvalues of the correlation are $0.224$, $1.077$, and $1.699$, so
this sign choice is positive definite. Naive Gaussian standard errors that
pretend the matrix is a sample correlation of 95 i.i.d.\ lines are about
$0.10$ for the two coefficients in the first equation and $0.084$ for
$0.590$. Campbell pooled seven environments, so those standard errors are not
his sampling variances. The point fit does not depend on $N$.

The scan was rerun with this $B$ and $\Psi$ and with the signed $G$ from the
glyph reading of equation~(7). Seed \texttt{20101029}, 100 replicates. Those
QTL coefficients were not estimated from the correlations, and on this
residual scale they are large: multi-trait power is 93\%, 95\%, and 99\%.
Single-trait power for QTL3 on $Y_1$ is 96\%. That is a consequence of the
scale mixture, not evidence that the 2003 correlations produce Mi's powers.
\ev{our own simulation, under an assumption we flag}
The bow check at $\rho=0.25$ (seed \texttt{20101030}) gives a diagonal
probability limit $0.198$ for the missing edge and a rejection rate $0.74$.
The unsigned-$G$ run of the same seed had multi-trait powers 100\%, 96\%, and
100\%. The sign change does not turn the calibration into a replication.

\section{Limitations}

The wheat fit is still blocked. It needs the 98 line identifiers, the ANOVA
residuals of grain yield, thousand-kernel weight, kernels per spike, and
spikes per square meter, the $\pm 1$ QTL and cofactor codes, the divisor-$n$
sample covariance, and the edge set as a matrix. Figure~5 draws arrows and
does not print that matrix. Campbell's Table~2 does not substitute for any
of those columns.

The Drton--Foygel--Sullivant criterion is not a theorem of this repository.
The delta method is not a Lean central limit theorem. The glyph reading of
equation~(7) meets Table~2's parameter block up to the two rounded QTL2
entries. It leaves the SEM powers for QTL2 and QTL3, and about eleven points
of the single-trait QTL3 power, short of Table~1. A second block tested
the choices the simulation paragraph does not pin down. None of them
closes the QTL2 gap. A $250$-line scan matches that one cell and misses
the first QTL. Those leftovers are not Monte Carlo error at $1000$ replicates. Twelve
generations of selfing, Haldane's map, the marker placement, and the
detection windows are our assumptions. The 2021 Hidalgo-Contreras wheat SEM
is a different experiment.

\section{Reproducibility}

The command list for every script is
\texttt{formalization/correlated-residuals/README.md}.
From \texttt{formalization/correlated-residuals/simulation},
\begin{verbatim}
python3 studies.py --reps 100 --bow-reps 400 --studies 1,2,3
\end{verbatim}
writes \texttt{study1.json}, \texttt{study2.json}, and
\texttt{study3.json}. The CSV tables are built from those files.
\begin{itemize}
\item Simulation commit: \texttt{5cd3c9f}.
\item Full hash: \sha{5cd3c9f6d2f6e38f25e32146d74d55f3778436e1}.
\item Branch: \texttt{cursor/l-f102-correlated-residuals-4c01}.
\item Lean 4.22.0 and Mathlib v4.22.0, built with \texttt{lake build}.
\item Theorem links are pinned to \sha{05d69ea141adf3d5ee936d3b48061bde241250dc}.
\item This PDF is \texttt{pdflatex paper.tex}. The ledger is \texttt{SOURCES.md}.
\item The gap grid is \texttt{python3 gaps.py} in the simulation directory. Seeds are \texttt{20101027} through \texttt{20101042}.
\end{itemize}

\appendix
\section{Central Lean proofs}
\label{app:lean}

The diagonal model reads the structural coefficients from the Schur complement
$\Omega=\Sigma_{yy}-\Sigma_{yx}\Sigma_{xx}^{-1}\Sigma_{yx}^{\mathsf T}$, then
recovers $\Gamma=(I-B)\Sigma_{yx}\Sigma_{xx}^{-1}$ and $\Phi=\Sigma_{xx}$.
The bow-free model drops the edge into the first trait from the third and
reads the residual covariance $c$ from the same complement. The collision
proofs exhibit, for an arbitrary exogenous coefficient, a second parameter
point with the same joint blocks. Consistency is continuous mapping of those
left inverses, composed with the strong law for the sample Gram matrix.
The listings below are the theorem bodies. Their lemmas live in the linked files.

\medskip
\noindent\texttt{recover\_sigma}, \texttt{Diagonal.lean}, lines 246--279.
\lstinputlisting[firstline=246,lastline=279]{../CorrelatedResiduals/Diagonal.lean}

\medskip
\noindent\texttt{recover\_bow}, \texttt{BowFree.lean}, lines 187--213.
\lstinputlisting[firstline=187,lastline=213]{../CorrelatedResiduals/BowFree.lean}

\medskip
\noindent\texttt{bow12\_same\_law} in \texttt{Collision.lean}, lines 114--128.
The siblings \texttt{bow13\_same\_law} and \texttt{bow23\_same\_law} repeat
the argument on the other two pairs, lines 130--160.
\lstinputlisting[firstline=114,lastline=128]{../CorrelatedResiduals/Collision.lean}

\medskip
\noindent\texttt{unitUpper\_factor\_unique}, \texttt{Triangular.lean}, lines 68--135.
\lstinputlisting[firstline=68,lastline=135]{../CorrelatedResiduals/Triangular.lean}

\medskip
\noindent\texttt{sigmaRec\_injective}, \texttt{Triangular.lean}, lines 298--317.
\lstinputlisting[firstline=298,lastline=317]{../CorrelatedResiduals/Triangular.lean}

\medskip
\noindent\texttt{plug\_in\_consistent\_ae}, \texttt{Law.lean}, lines 174--199.
\lstinputlisting[firstline=174,lastline=199]{../CorrelatedResiduals/Law.lean}

\medskip
\noindent\texttt{plug\_in\_bow\_consistent\_ae}, \texttt{BowConsistency.lean}, lines 270--304.
\lstinputlisting[firstline=270,lastline=304]{../CorrelatedResiduals/BowConsistency.lean}

\section{Artifacts}
\label{app:artifacts}

Each row is one claim in this note, the file that holds it, and the command
that regenerates it. Commands are from the repository root unless the row
says otherwise. Commits are on
\texttt{cursor/l-f102-correlated-residuals-4c01}. A short \texttt{--reps}
run overwrites the committed JSON; the commands below are the runs that
match the cited files.

\begin{small}
\begin{longtable}{@{}>{\raggedright\arraybackslash}p{0.30\textwidth}>{\raggedright\arraybackslash}p{0.32\textwidth}>{\raggedright\arraybackslash}p{0.30\textwidth}@{}}
\toprule
Claim & File and commit & Command \\
\midrule
\endhead
Glyph boxes for equation~(7): equals ends at $145.9$, parenthesis at $148.1$, four unary minuses &
\sha{simulation/out/qtl1/glyphs.json}, \sha{b4e8a3c594f12f419560d5cb2a569b649ae93a55} &
\sha{python3 formalization/correlated-residuals/simulation/measure_glyphs.py --pdf local-extraction/literature/unl-sem/mi-qtl2010.pdf --check formalization/correlated-residuals/simulation/out/qtl1/glyphs.json} \\
\addlinespace
Total effects, including $3/4$ and $-7/16$ &
\sha{simulation/out/table2_totals.json} &
\sha{python3 formalization/correlated-residuals/simulation/table2_totals.py} \\
\addlinespace
Baseline SEM powers $50.8$, $53.5$, $86.1$ &
\sha{simulation/out/assumptions/baseline.json}, \sha{82ff59d6bcfb02c273aaf1f83195e0d120fe0915} &
\sha{python3 assumptions.py --configs baseline --reps 1000} in \texttt{simulation/} \\
\addlinespace
Gap grid, including the $13\%$ and $86.2\%$ single-trait cells &
\sha{simulation/out/gaps.json}, \sha{gaps_table.csv}, \sha{c2769042e51806a77b2b70237fb4c014552e78fb} &
\sha{python3 gaps.py} in \texttt{simulation/} \\
\addlinespace
Closest triple $65.2$, $73.3$, $99.3$ &
\sha{simulation/out/qtl1/qtl1_24_n250.json}, \sha{b4e8a3c594f12f419560d5cb2a569b649ae93a55} &
\sha{python3 qtl1_search.py --configs qtl1_24_n250 --reps 1000} in \texttt{simulation/} \\
\addlinespace
Product $1.37\times 10^{-12}$ &
\sha{simulation/out/printed_probability.json}, from the baseline file above &
\sha{python3 formalization/correlated-residuals/simulation/printed_probability.py} \\
\addlinespace
Studies 1--3, bow Monte Carlo, Campbell scale check &
\sha{simulation/out/study1.json}, \sha{study2.json}, \sha{study3.json}, \sha{5cd3c9f6d2f6e38f25e32146d74d55f3778436e1} &
\sha{python3 studies.py --reps 100 --bow-reps 400 --studies 1,2,3} in \texttt{simulation/} \\
\addlinespace
Identification algebra at the earlier sign convention &
\sha{witnesses.json}, written by \sha{check_witness.py} &
\sha{python3 check_witness.py} in \sha{formalization/correlated-residuals/} \\
\addlinespace
Delta-method numeric check &
\sha{delta_check.json} &
\sha{python3 delta_method.py} in \sha{formalization/correlated-residuals/} \\
\addlinespace
Bow-free plug-in numeric check &
\sha{fit_bow.json} &
\sha{python3 fit_bow.py} in \sha{formalization/correlated-residuals/} \\
\addlinespace
Lean theorems, no \texttt{sorry}. Axioms \texttt{propext}, \texttt{Classical.choice}, \texttt{Quot.sound} &
\sha{CorrelatedResiduals/*.lean}, listings pinned to \sha{05d69ea141adf3d5ee936d3b48061bde241250dc} &
\sha{lake build} in \sha{formalization/correlated-residuals/} \\
\addlinespace
This PDF &
\sha{paper/paper.pdf} &
\sha{pdflatex paper.tex} in \texttt{paper/} \\
\bottomrule
\end{longtable}
\end{small}

\begin{thebibliography}{9}
\bibitem{brito2002}
C.~Brito and J.~Pearl.
A new identification condition for recursive models with correlated errors.
\emph{Structural Equation Modeling} 9 (2002), 459--474.

\bibitem{campbell2003}
B.~T.~Campbell, P.~S.~Baenziger, K.~S.~Gill, K.~M.~Eskridge, H.~Budak,
M.~Erayman, and Y.~Yen.
Identification of QTLs and environmental interactions associated with
agronomic traits on chromosome 3A of wheat.
\emph{Crop Science} 43 (2003), 1493--1505.

\bibitem{drton2011}
M.~Drton, R.~Foygel, and S.~Sullivant.
Global identifiability of linear structural equation models.
\emph{Annals of Statistics} 39 (2011), 865--886.

\bibitem{foygel2012}
R.~Foygel, J.~Draisma, and M.~Drton.
Half-trek criterion for generic identifiability of linear structural
equation models.
\emph{Annals of Statistics} 40 (2012), 1682--1713.

\bibitem{hidalgo2021}
J.~V.~Hidalgo-Contreras, J.~Salinas-Ruiz, K.~M.~Eskridge, and
S.~P.~Baenziger.
Incorporating molecular markers and causal structure among traits using a
Smith--Hazel index and structural equation models.
\emph{Agronomy} 11 (2021), 1953.

\bibitem{li2006}
R.~Li, S.-W.~Tsaih, K.~Shockley, I.~M.~Stylianou, J.~Wergedal, B.~Paigen,
and G.~A.~Churchill.
Structural model analysis of multiple quantitative traits.
\emph{PLoS Genetics} 2 (2006), e114.

\bibitem{mi2008}
X.~Mi, K.~Eskridge, and D.~Wang.
Multi-trait QTL mapping using a structural equation model.
Conference on Applied Statistics in Agriculture (2008).
DOI \texttt{10.4148/2475-7772.1095}.

\bibitem{mi2009}
X.~Mi.
Genetic applications using structural equation modeling.
Doctoral dissertation, University of Nebraska--Lincoln, 2009.
ProQuest AAI3350451. Not consulted. The full text was not obtained.

\bibitem{mi2010}
X.~Mi, K.~Eskridge, D.~Wang, P.~S.~Baenziger, B.~T.~Campbell, K.~S.~Gill,
I.~Dweikat, and J.~Bovaird.
Regression-based multi-trait QTL mapping using a structural equation model.
\emph{Statistical Applications in Genetics and Molecular Biology} 9 (2010),
article~38.

\bibitem{mi2010gr}
X.~Mi, K.~Eskridge, D.~Wang, P.~S.~Baenziger, B.~T.~Campbell, K.~S.~Gill,
and I.~Dweikat.
Bayesian mixture structural equation modelling in multiple-trait QTL mapping.
\emph{Genetics Research} 92 (2010), 239--250.

\bibitem{valente2010}
B.~D.~Valente, G.~J.~M.~Rosa, G.~de los Campos, D.~Gianola, and M.~A.~Silva.
Searching for Recursive Causal Structures in Multivariate Quantitative
Genetics Mixed Models.
\emph{Genetics} 185 (2010), 633--644.
\end{thebibliography}

\end{document}
